% ============================================================================
% OpenDSSSim P1 Volt/VAR + DER Controller Science
% STATUS: DAY-2 CONTROLLER MATHEMATICS RATIFIED
% Date: 2026-10-10
%
% This document is complete at the controller-mathematics level for P1.
% Remaining experiment/configuration items before final campaign freeze:
%   (i) exact IEEE-123 device placement/ratings;
%   (ii) one fixed numerical common-cost weight set after scale verification;
%   (iii) RL training-budget/hyperparameter qualification.
%
% No generalized controller implementation may hard-code IEEE-123 object names,
% device counts, phase counts, or experiment default values.
% ============================================================================

\section{P1 Volt/VAR and DER Coordination}
\label{sec:p1_control_science}

\subsection{Scope, sources, and design principle}
\label{subsec:p1_scope}

The P1 benchmark defines one physical Volt/VAR coordination problem and three
controller paradigms:
\begin{enumerate}
    \item a deterministic local rule-based controller (RBC);
    \item a centralized LinDistFlow mixed-integer quadratic model predictive controller
          (MIQP-MPC);
    \item an MPC-inspired, model-free branching dueling double deep Q-network (BDQ).
\end{enumerate}
All three controllers operate the same OpenDSSSim plant and the same selected physical
actuators.  Their distinction lies in the decision mechanism, not in the plant.

The physical VVC/Gym framing is grounded in PowerGym~\cite{fan2022powergym} and the
graph-policy VVC study of Lee et al.~\cite{lee2022graph}.  Forecast-aware unbalanced
predictive control with taps, capacitors, and inverter-based DER is informed by
Erazo-Caicedo et al.~\cite{erazo2022model}.  The exact radial DistFlow starting point
follows Baran and Wu~\cite{baran1989sizing}; the multiphase linearized extension is
consistent with established linear approximations for unbalanced radial networks
\cite{gan2014multiphase}.  The high-dimensional discrete RL architecture follows
Tavakoli et al.~\cite{tavakoli2018branching}.  The OpenDSS native control semantics used
to motivate the RBC are documented by the OpenDSS project
\cite{opendss_invcontrol,opendss_capcontrol,opendss_regcontrol}.

The controller mathematics is cardinality-agnostic.  The selected IEEE-123 experiment
may instantiate a particular number of devices, but the controller equations and source
must not assume those counts.

\subsection{Network and device sets}
\label{subsec:p1_sets}

Let the radial feeder be a directed graph
\[
\mathcal G=(\mathcal N,\mathcal E),
\]
oriented away from the source/root bus $0$.  Let $\Phi_i\subseteq\{a,b,c\}$ denote the
phases present at bus $i$, and let
\[
\mathcal B_{\Phi}
=
\{(i,\phi):i\in\mathcal N,\ \phi\in\Phi_i\}
\]
be the active bus--phase set.

The selected controllable device-channel sets are
\[
\mathcal R,\qquad
\mathcal C,\qquad
\mathcal P,
\]
for regulator/tap, capacitor, and PV-inverter reactive-power channels, respectively, with
\[
N_R=|\mathcal R|,\qquad
N_C=|\mathcal C|,\qquad
N_P=|\mathcal P|.
\]
No equation below assumes fixed $N_R$, $N_C$, or $N_P$.  The PESGM IEEE-123 experiment
is expected to use ten channels of each type, but this is an experiment configuration,
not a controller assumption.

\subsection{Common time and exogenous information}
\label{subsec:p1_time_info}

For the PESGM experiment the source-backed control interval and predictive horizon are
\[
\Delta t = 5~\mathrm{min},
\qquad
H=12,
\qquad
H\Delta t = 1~\mathrm{h},
\]
following the predictive timing used by Erazo-Caicedo et al.~\cite{erazo2022model}.
The generalized implementations keep both $\Delta t$ and $H$ configurable.

At time $t$, the exogenous electrical input is
\[
d_t
=
\begin{bmatrix}
(\mathbf p_t^L)^\top &
(\mathbf q_t^L)^\top &
(\mathbf p_t^{PV,\mathrm{av}})^\top
\end{bmatrix}^{\top}.
\]
PV active power follows available generation,
\[
\mathbf p_t^{PV}
=
\mathbf p_t^{PV,\mathrm{av}},
\]
and is not a P1 control variable.

MPC and RL receive the same authorized forecast tensor
\[
\widehat{\mathbf D}_t
=
\left\{
\hat d_{t+k|t}
\right\}_{k=0}^{H-1},
\]
where
\[
\hat d_{t+k|t}
=
\begin{bmatrix}
(\hat{\mathbf p}_{t+k|t}^{L})^\top &
(\hat{\mathbf q}_{t+k|t}^{L})^\top &
(\hat{\mathbf p}_{t+k|t}^{PV,\mathrm{av}})^\top
\end{bmatrix}^{\top}.
\]
Forecast access by RL is an explicit benchmark adaptation for fair information access; the
Lee et al. source policy itself does not require this forecast tensor.

\subsection{Common observation and requested action}
\label{subsec:p1_obs_action}

Let $x_t$ denote the full physical plant state and
\[
o_t=h(x_t)
\]
the authorized controller observation.  The P1 observation contains the available
bus--phase voltage magnitudes, realized regulator states, realized capacitor states, PV
active/reactive operating points, and authorized current load/PV quantities.  MPC and RL
may additionally use $\widehat{\mathbf D}_t$.

Every paradigm emits the same physical requested-action structure,
\[
a_t^{\mathrm{req}}
=
\begin{bmatrix}
\Delta\boldsymbol\tau_t^\top &
(\mathbf c_t^{\mathrm{cmd}})^\top &
(\mathbf q_t^{PV,\mathrm{cmd}})^\top
\end{bmatrix}^{\top},
\]
with
\[
\Delta\tau_{r,t}\in\{-1,0,+1\},
\qquad r\in\mathcal R,
\]
\[
c_{c,t}^{\mathrm{cmd}}\in\{0,1\},
\qquad c\in\mathcal C.
\]
PV reactive-power commands satisfy the instantaneous inverter capability defined below.

OpenDSSSim owns
\[
a_t^{\mathrm{req}}
\longrightarrow
a_t^{\mathrm{app}}
\longrightarrow
a_t^{\mathrm{real}}
\]
and the authoritative plant transition
\[
x_{t+1}
=
F_{\mathrm{DSS}}
\!\left(
x_t,a_t^{\mathrm{app}},d_t
\right).
\]
Controllers never directly mutate AltDSS.  Each control interval has exactly one
\texttt{env.step(...)} call containing the complete requested action vector.

\subsection{Common physical device constraints}
\label{subsec:p1_device_constraints}

For inverter $p\in\mathcal P$, with nameplate apparent-power rating
$S_p^{\mathrm{rated}}$,
\[
\left(P_{p,t}^{PV,\mathrm{av}}\right)^2
+
\left(q_{p,t}^{PV}\right)^2
\le
\left(S_p^{\mathrm{rated}}\right)^2.
\]
Since active PV power is fixed at forecast/realized availability, define
\[
Q_{p,t}^{\max}
=
\sqrt{
\left(S_p^{\mathrm{rated}}\right)^2
-
\left(P_{p,t}^{PV,\mathrm{av}}\right)^2
},
\]
giving
\[
-Q_{p,t}^{\max}
\le
q_{p,t}^{PV}
\le
Q_{p,t}^{\max}.
\]

The desired operating voltage band is
\[
\underline V=0.95,
\qquad
\overline V=1.05
\quad\mathrm{p.u.},
\]
consistent with the P1 source literature
\cite{fan2022powergym,erazo2022model,lee2022graph}.

\subsection{Common stage-cost semantics}
\label{subsec:p1_common_cost}

Let
\[
e_{i,\phi,t}^{V,+}
=
\left(|V_{i,\phi,t}|-\overline V\right)^+,
\qquad
e_{i,\phi,t}^{V,-}
=
\left(\underline V-|V_{i,\phi,t}|\right)^+,
\]
where $(z)^+=\max(z,0)$.  A cardinality-normalized voltage-violation metric is
\[
J_t^V
=
\frac{1}{|\mathcal B_\Phi|}
\sum_{(i,\phi)\in\mathcal B_\Phi}
\left[
\left(e_{i,\phi,t}^{V,+}\right)^2
+
\left(e_{i,\phi,t}^{V,-}\right)^2
\right].
\]

Mechanical switching effort is
\[
J_t^{\mathrm{sw}}
=
\frac{1}{N_R+N_C}
\left[
\sum_{r\in\mathcal R}
\left|\Delta\tau_{r,t}\right|
+
\sum_{c\in\mathcal C}
\left|c_{c,t}-c_{c,t-1}\right|
\right].
\]

Normalized inverter reactive-power effort is
\[
J_t^Q
=
\frac{1}{N_P}
\sum_{p\in\mathcal P}
\left(
\frac{q_{p,t}^{PV}}
{\max(Q_{p,t}^{\max},\epsilon_Q)}
\right)^2,
\]
where $\epsilon_Q>0$ prevents division by zero when no reactive headroom remains.

The common stage cost is
\[
\boxed{
\ell_t
=
w_V J_t^V
+
w_{\mathrm{sw}}J_t^{\mathrm{sw}}
+
w_QJ_t^Q
}
\]
with fixed positive weights satisfying the intended priority
\[
w_V\gg w_{\mathrm{sw}}>w_Q.
\]
The project uses one fixed basic weight assignment after scale verification and does not
perform a Pareto study or select weights using test-set controller performance.

% ============================================================================
\subsection{Rule-based controller}
\label{subsec:p1_rbc}

The RBC is model-free and defined entirely by local engineering logic,
\[
a_t^{\mathrm{RBC}}
=
\mu_{\mathrm{RBC}}(o_t).
\]

\subsubsection{Configurable local monitor reduction}

For each device channel $j$, let $\mathbf V_{j,t}^{\mathrm{mon}}$ contain the configured
voltage channels associated with that device.  A configurable reduction operator
$\rho_j$ forms the scalar local control voltage
\[
V_{j,t}^{\mathrm{mon}}
=
\rho_j\!\left(
\mathbf V_{j,t}^{\mathrm{mon}}
\right).
\]
The generalized RBC may support configured phase selection and reductions such as
\texttt{AVG}, \texttt{MIN}, or \texttt{MAX} as applicable.  For the PESGM experiment,
all multiphase P1 RBC monitors use
\[
\boxed{
\rho_j=\mathrm{AVG}.
}
\]
For a single-phase channel this reduces trivially to that phase voltage.

\subsubsection{PV Volt/VAR rule}

For inverter $p$ define the normalized reactive command
\[
q_{p,t}^{PV,\mathrm{req}}
=
Q_{p,t}^{\max}\,
g_p\!\left(
V_{p,t}^{\mathrm{mon}}
\right),
\]
where the generic breakpoint set is configurable.  The PESGM default curve is
\[
g_p(V)
=
\begin{cases}
+1,
&
V\le 0.95,
\\[2mm]
\dfrac{0.99-V}{0.04},
&
0.95<V<0.99,
\\[3mm]
0,
&
0.99\le V\le1.01,
\\[2mm]
\dfrac{1.01-V}{0.04},
&
1.01<V<1.05,
\\[3mm]
-1,
&
V\ge1.05.
\end{cases}
\]
Thus undervoltage requests reactive injection, overvoltage requests reactive absorption,
and the request is always scaled by the currently available inverter headroom.
The Volt/VAR principle is source-backed by native inverter-control practice
\cite{opendss_invcontrol}; the exact four breakpoints above are an OpenDSSSim experiment
adaptation.

\subsubsection{Capacitor hysteresis rule}

For capacitor channel $c$,
\[
c_{c,t}^{\mathrm{req}}
=
\begin{cases}
1,
&
V_{c,t}^{\mathrm{mon}}
\le V_c^{\mathrm{on}},
\\
0,
&
V_{c,t}^{\mathrm{mon}}
\ge V_c^{\mathrm{off}},
\\
c_{c,t-1}^{\mathrm{real}},
&
V_c^{\mathrm{on}}
<
V_{c,t}^{\mathrm{mon}}
<
V_c^{\mathrm{off}}.
\end{cases}
\]
The generic thresholds are configurable.  The PESGM defaults are
\[
\boxed{
V_c^{\mathrm{on}}=0.97,
\qquad
V_c^{\mathrm{off}}=1.03.
}
\]
Voltage-hysteresis switching follows conventional capacitor-control semantics
\cite{opendss_capcontrol}; the numerical thresholds are an OpenDSSSim experiment
adaptation.

\subsubsection{Regulator deadband rule}

For regulator channel $r$, let $V_r^{\mathrm{ref}}$ and $B_r$ denote the configurable
per-unit reference and total deadband width.  Then
\[
\Delta\tau_{r,t}^{\mathrm{req}}
=
\begin{cases}
+1,
&
V_{r,t}^{\mathrm{mon}}
<
V_r^{\mathrm{ref}}-\dfrac{B_r}{2},
\\[2mm]
0,
&
\left|
V_{r,t}^{\mathrm{mon}}-V_r^{\mathrm{ref}}
\right|
\le
\dfrac{B_r}{2},
\\[2mm]
-1,
&
V_{r,t}^{\mathrm{mon}}
>
V_r^{\mathrm{ref}}+\dfrac{B_r}{2}.
\end{cases}
\]
The sign is semantic: $+1$ requests one step that raises the controlled-side voltage and
$-1$ requests one step that lowers it.  The PESGM defaults are
\[
\boxed{
V_r^{\mathrm{ref}}=1.0,
\qquad
B_r=0.025,
}
\]
giving the deadband
\[
[0.9875,\,1.0125]\ \mathrm{p.u.}
\]
These values are compatible with conventional RegControl reference/band concepts
\cite{opendss_regcontrol}; the generic controller keeps them configurable.

\subsubsection{Simultaneous local execution}

There is no sequential PV--capacitor--regulator plant solve and no P1 controller delay.
All local rules evaluate from the same observation,
\[
a_t^{\mathrm{RBC}}
=
\begin{bmatrix}
\{\Delta\tau_{r,t}^{\mathrm{req}}\}_{r\in\mathcal R}\\
\{c_{c,t}^{\mathrm{req}}\}_{c\in\mathcal C}\\
\{q_{p,t}^{PV,\mathrm{req}}\}_{p\in\mathcal P}
\end{bmatrix},
\]
and the complete vector is applied through one
\[
\texttt{env.step}(a_t^{\mathrm{RBC}}).
\]

% ============================================================================
\subsection{Exact radial DistFlow starting point}
\label{subsec:p1_distflow}

Consider a single-phase branch $i\rightarrow j$ with series impedance
\[
z_{ij}=r_{ij}+\mathrm j x_{ij}.
\]
Ohm's law is
\[
V_j
=
V_i-z_{ij}I_{ij}.
\]
Define squared voltage and squared current magnitude
\[
v_i=|V_i|^2,
\qquad
\ell_{ij}=|I_{ij}|^2,
\]
and sending-end complex power
\[
S_{ij}
=
P_{ij}+\mathrm jQ_{ij}
=
V_iI_{ij}^{*}.
\]
Taking the squared magnitude of Ohm's law yields
\[
\boxed{
v_j
=
v_i
-
2\left(
r_{ij}P_{ij}
+
x_{ij}Q_{ij}
\right)
+
\left(
r_{ij}^{2}+x_{ij}^{2}
\right)\ell_{ij}.
}
\]
The magnitude relation $|S_{ij}|^2=|V_i|^2|I_{ij}|^2$ gives
\[
\boxed{
v_i\ell_{ij}
=
P_{ij}^2+Q_{ij}^2.
}
\]

Let $\mathcal C(j)$ denote the child buses downstream of $j$.  Including branch losses,
the exact radial branch balances are
\[
\boxed{
P_{ij}
=
p_j
+
\sum_{m\in\mathcal C(j)}
P_{jm}
+
r_{ij}\ell_{ij},
}
\]
\[
\boxed{
Q_{ij}
=
q_j
+
\sum_{m\in\mathcal C(j)}
Q_{jm}
+
x_{ij}\ell_{ij}.
}
\]
These equations form the standard DistFlow model
\cite{baran1989sizing}.

\subsection{LinDistFlow approximation}
\label{subsec:p1_lindist}

The internal MPC model neglects the terms proportional to $\ell_{ij}$:
\[
r_{ij}\ell_{ij},
\qquad
x_{ij}\ell_{ij},
\qquad
(r_{ij}^2+x_{ij}^2)\ell_{ij}.
\]
The resulting lossless radial LinDistFlow equations are
\[
\boxed{
P_{ij}
=
p_j
+
\sum_{m\in\mathcal C(j)}
P_{jm},
}
\]
\[
\boxed{
Q_{ij}
=
q_j
+
\sum_{m\in\mathcal C(j)}
Q_{jm},
}
\]
\[
\boxed{
v_j
=
v_i
-
2(r_{ij}P_{ij}+x_{ij}Q_{ij}).
}
\]
The nonlinear current equation is then no longer required in the internal prediction
model.

\subsection{Embedding PV and capacitor controls in LinDistFlow}
\label{subsec:p1_lindist_devices}

With the convention that positive $p_j$ and $q_j$ denote net consumption,
\[
p_{j,k}
=
\hat p_{j,k}^{L}
-
\hat p_{j,k}^{PV,\mathrm{av}},
\]
and
\[
q_{j,k}
=
\hat q_{j,k}^{L}
-
q_{j,k}^{PV}
-
q_{j,k}^{cap}.
\]

A switched capacitor with rated reactive injection $\bar q_c^{cap}$ is modeled internally
as
\[
q_{c,k}^{cap}
=
\bar q_c^{cap}c_{c,k},
\qquad
c_{c,k}\in\{0,1\}.
\]
This constant-kvar approximation is deliberate: native voltage-dependent capacitor
behavior remains in the OpenDSS plant, while the MPC prediction remains affine.

Therefore PV and capacitor control enter the network through reactive balance,
\[
Q_{ij,k}
=
\hat q_{j,k}^{L}
-
q_{j,k}^{PV}
-
\bar q_j^{cap}c_{j,k}
+
\sum_{m\in\mathcal C(j)}
Q_{jm,k},
\]
with the obvious summation/mapping generalization when multiple devices connect to the
same bus/phase.  Their voltage effect follows the chain
\[
q^{PV},c
\longrightarrow
q_j
\longrightarrow
Q_{ij}
\longrightarrow
v_j.
\]

\subsection{Embedding regulator taps in LinDistFlow}
\label{subsec:p1_reg_lindist}

For an ideal regulator/transformer with voltage ratio $a_{ij}$,
\[
\widetilde V_i
=
a_{ij}V_i
\]
and hence
\[
\widetilde v_i
=
a_{ij}^{2}v_i.
\]
With series impedance after the ratio transformation, the corresponding lossless
single-phase voltage relation is conceptually
\[
v_j
\simeq
a_{ij}^{2}v_i
-
2(r_{ij}P_{ij}+x_{ij}Q_{ij}).
\]

Let tap position $\tau_r$ parameterize the ratio around a nominal value:
\[
a_r(\tau_r)
=
\bar a_r
+
\Delta a_r
(\tau_r-\bar\tau_r).
\]
For
\[
f(a,v)=a^2v,
\]
the first-order expansion about $(\bar a,\bar v)$ is
\[
a^2v
\approx
\bar a^2v
+
2\bar a\bar v(a-\bar a).
\]
Around nominal $\bar a\simeq1$ and $\bar v\simeq1$,
\[
a^2v
\approx
v
+
2\Delta a_r
(\tau_r-\bar\tau_r).
\]
Define
\[
\kappa_r
=
2\bar a_r\bar v_r\Delta a_r
\]
(or the circuit/model-service-calibrated equivalent linear coefficient), giving the
affine tap contribution
\[
\boxed{
v_j
=
v_i
-
2(r_{ij}P_{ij}+x_{ij}Q_{ij})
+
\kappa_r(\tau_r-\bar\tau_r).
}
\]

Tap position is a persistent actuator state:
\[
\boxed{
\tau_{r,k}
=
\tau_{r,k-1}
+
\Delta\tau_{r,k},
}
\]
with
\[
\Delta\tau_{r,k}\in\{-1,0,+1\},
\qquad
\tau_r^{\min}
\le
\tau_{r,k}
\le
\tau_r^{\max}.
\]
Thus $\Delta\tau$ is the movement command, while $\tau$ is the absolute predicted tap
position used by the voltage equation.

\subsection{Multiphase unbalanced LinDistFlow}
\label{subsec:p1_multiphase}

For a branch $(i,j)$, collect only the physically present phase quantities in
\[
\mathbf P_{ij,k},
\qquad
\mathbf Q_{ij,k},
\qquad
\mathbf v_{i,k}.
\]
Phase-wise branch balance is
\[
\mathbf P_{ij,k}
=
\mathbf p_{j,k}
+
\sum_{m\in\mathcal C(j)}
\mathbf P_{jm,k},
\]
\[
\mathbf Q_{ij,k}
=
\mathbf q_{j,k}
+
\sum_{m\in\mathcal C(j)}
\mathbf Q_{jm,k}.
\]
A phase-aware linear voltage relation is written
\[
\boxed{
\mathbf v_{j,k}
=
\mathbf v_{i,k}
-
\mathbf R_{ij}\mathbf P_{ij,k}
-
\mathbf X_{ij}\mathbf Q_{ij,k}
+
\mathbf T_{ij}
(\boldsymbol\tau_{ij,k}-\bar{\boldsymbol\tau}_{ij}).
}
\]
The matrices $\mathbf R_{ij}$ and $\mathbf X_{ij}$ retain same-phase and cross-phase
coupling; $\mathbf T_{ij}$ maps the available regulator tap channels to affected phase
voltages.  This form is consistent with established multiphase linear approximations for
radial distribution systems~\cite{gan2014multiphase}.

Device-to-phase incidence matrices map controller channels to network injections:
\[
\mathbf q_{k}^{PV,\mathrm{bus}}
=
\mathbf M_{PV}\mathbf q_k^{PV},
\]
\[
\mathbf q_{k}^{cap,\mathrm{bus}}
=
\mathbf M_C
\operatorname{diag}(\bar{\mathbf q}^{cap})
\mathbf c_k.
\]
The active net-injection vector is
\[
\mathbf p_k
=
\hat{\mathbf p}_k^L
-
\hat{\mathbf p}_k^{PV,\mathrm{av}},
\]
and the reactive net-consumption vector is
\[
\mathbf q_k
=
\hat{\mathbf q}_k^L
-
\mathbf M_{PV}\mathbf q_k^{PV}
-
\mathbf M_C
\operatorname{diag}(\bar{\mathbf q}^{cap})
\mathbf c_k.
\]

After eliminating branch-flow variables, the model can be represented in compact affine
form
\[
\boxed{
\mathbf v_k
=
\mathbf b_k
+
\mathbf G_{PV}\mathbf q_k^{PV}
+
\mathbf G_C\mathbf c_k
+
\mathbf G_R\boldsymbol\tau_k,
}
\]
where $\mathbf b_k$ contains the source-voltage and forecasted uncontrollable-injection
contribution.  The branch-level equations remain the authoritative explanation of the
physics; the compact form is the implementation convenience.

% ============================================================================
\subsection{MIQP-MPC formulation}
\label{subsec:p1_mpc}

At control time $t$, let
\[
\mathcal K=\{0,\ldots,H-1\}.
\]
The continuous decision variables include
\[
\mathbf P_{ij,k},
\quad
\mathbf Q_{ij,k},
\quad
\mathbf v_{i,k},
\quad
\mathbf q_k^{PV},
\quad
\boldsymbol\xi_k^{+},
\quad
\boldsymbol\xi_k^{-},
\]
and nonnegative switching auxiliaries.  The discrete variables include
\[
\mathbf c_k\in\{0,1\}^{N_C},
\]
\[
\boldsymbol\tau_k\in\mathbb Z^{N_R},
\qquad
\Delta\boldsymbol\tau_k\in\{-1,0,+1\}^{N_R}.
\]

The voltage-band constraints use nonnegative slacks
\[
\underline V^2-\xi_{i,\phi,k}^{-}
\le
v_{i,\phi,k}
\le
\overline V^2+\xi_{i,\phi,k}^{+},
\]
\[
\xi_{i,\phi,k}^{+}\ge0,
\qquad
\xi_{i,\phi,k}^{-}\ge0,
\]
to prevent artificial optimization infeasibility when the forecasted network cannot
satisfy the desired voltage band.

PV capability is
\[
-Q_{p,k}^{\max}
\le
q_{p,k}^{PV}
\le
Q_{p,k}^{\max},
\]
with
\[
Q_{p,k}^{\max}
=
\sqrt{
(S_p^{\mathrm{rated}})^2
-
(\hat P_{p,k}^{PV,\mathrm{av}})^2
}
\]
precomputed for each horizon step.

Tap dynamics are
\[
\tau_{r,k}
=
\tau_{r,k-1}
+
\Delta\tau_{r,k},
\]
with the first-step predecessor fixed by the currently realized plant tap,
\[
\tau_{r,-1|t}
=
\tau_{r,t}^{\mathrm{real}}.
\]
Tap bounds are
\[
\tau_r^{\min}
\le
\tau_{r,k}
\le
\tau_r^{\max}.
\]

For regulator movement introduce $s_{r,k}^{R}\ge0$ satisfying
\[
s_{r,k}^{R}
\ge
\Delta\tau_{r,k},
\qquad
s_{r,k}^{R}
\ge
-\Delta\tau_{r,k},
\]
so that minimization realizes
\[
s_{r,k}^{R}=|\Delta\tau_{r,k}|.
\]
For capacitor switching introduce $s_{c,k}^{C}\ge0$ satisfying
\[
s_{c,k}^{C}
\ge
c_{c,k}-c_{c,k-1},
\]
\[
s_{c,k}^{C}
\ge
-(c_{c,k}-c_{c,k-1}),
\]
with
\[
c_{c,-1|t}
=
c_{c,t}^{\mathrm{real}}.
\]

The MPC solves
\[
\boxed{
\begin{aligned}
\min\quad
&
\sum_{k\in\mathcal K}
\Bigg[
w_V
\frac{1}{|\mathcal B_\Phi|}
\sum_{(i,\phi)\in\mathcal B_\Phi}
\left(
(\xi_{i,\phi,k}^{+})^2
+
(\xi_{i,\phi,k}^{-})^2
\right)
\\
&\qquad
+
w_{\mathrm{sw}}
\frac{
\sum_{r\in\mathcal R}s_{r,k}^{R}
+
\sum_{c\in\mathcal C}s_{c,k}^{C}
}{
N_R+N_C
}
\\
&\qquad
+
w_Q
\frac{1}{N_P}
\sum_{p\in\mathcal P}
\left(
\frac{q_{p,k}^{PV}}
{\max(Q_{p,k}^{\max},\epsilon_Q)}
\right)^2
\Bigg]
\end{aligned}
}
\]
subject to the complete multiphase LinDistFlow balance, voltage, source/slack, mapping,
tap, capacitor, PV capability, and variable-domain constraints above.

All network relations are affine and the objective is convex quadratic in the continuous
variables.  With integer/binary tap and capacitor decisions, this is a
Gurobi-compatible MIQP.

The receding-horizon controller applies only the first optimized action,
\[
\boxed{
a_t^{\mathrm{MPC}}
=
u_{t|t}^{\star},
}
\]
through the common OpenDSSSim \texttt{env.step(...)} boundary, receives the realized plant
state, shifts the horizon, and resolves.

% ============================================================================
\subsection{MPC-inspired model-free branching DQN}
\label{subsec:p1_rl}

The RL controller uses the same physical problem semantics but no LinDistFlow or other
transition model.  Its policy input is
\[
s_t^{RL}
=
\begin{bmatrix}
o_t\\
\widehat{\mathbf D}_t
\end{bmatrix},
\]
after fixed normalization of continuous features.

\subsubsection{Branch action sets}

The number of action branches is
\[
N_A
=
N_R+N_C+N_P.
\]
For regulator branch $r$,
\[
\mathcal A_r^R
=
\{-1,0,+1\}.
\]
For capacitor branch $c$,
\[
\mathcal A_c^C
=
\{0,1\}.
\]
For PV branch $p$,
\[
\boxed{
\mathcal A_p^{PV}
=
\{-1,-0.5,0,+0.5,+1\}.
}
\]
If $\alpha_{p,t}\in\mathcal A_p^{PV}$ is selected,
\[
\boxed{
q_{p,t}^{PV,\mathrm{req}}
=
\alpha_{p,t}Q_{p,t}^{\max}.
}
\]
Thus the logical action remains normalized while its physical kvar value follows the
time-varying inverter capability.

A flat DQN would require the Cartesian action space
\[
3^{N_R}2^{N_C}5^{N_P},
\]
which grows combinatorially.  The branching architecture of
Tavakoli et al.~\cite{tavakoli2018branching} instead uses a shared decision module and
one output branch per action dimension.

\subsubsection{Shared encoder and dueling branch values}

Let the shared encoder produce
\[
h_t
=
f_\theta(s_t^{RL}).
\]
The common value stream produces $V(s_t)$ and branch $d$ produces advantages
$A_d(s_t,a_d)$.  The branch action value is
\[
\boxed{
Q_d(s_t,a_d)
=
V(s_t)
+
A_d(s_t,a_d)
-
\frac{1}{|\mathcal A_d|}
\sum_{a_d'\in\mathcal A_d}
A_d(s_t,a_d').
}
\]
The greedy branch action is
\[
a_{d,t}^{\star}
=
\arg\max_{a_d\in\mathcal A_d}
Q_d(s_t,a_d).
\]
All branchwise choices are concatenated to form one simultaneous full requested action,
\[
a_t^{RL}
=
[
a_{1,t}^{\star},
\ldots,
a_{N_A,t}^{\star}
].
\]
The action is then applied by one
\[
\texttt{env.step}(a_t^{RL}).
\]

\subsubsection{Reward}

The RL reward is the negative common physical stage cost,
\[
\boxed{
r_t=-\ell_t.
}
\]
Hence the learned controller is MPC-inspired through the shared physical objective,
device semantics, bounds, and forecast information, while remaining model-free because
its transition samples come exclusively from the OpenDSSSim environment.

\subsubsection{Branching double-DQN target and loss}

Let $\theta$ denote the online network and $\theta^{-}$ the target network.  The online
network selects the next greedy action in branch $d$,
\[
a_{d,t+1}^{\star}
=
\arg\max_{a_d'}
Q_d(s_{t+1},a_d';\theta),
\]
and the target network evaluates it.  Using the mean branch target,
\[
\boxed{
y_t
=
r_t
+
\gamma(1-\delta_t)
\frac{1}{N_A}
\sum_{d=1}^{N_A}
Q_d
\left(
s_{t+1},
a_{d,t+1}^{\star};
\theta^{-}
\right),
}
\]
where $\delta_t$ indicates a true terminal transition.

For a replayed transition, the training loss is
\[
\boxed{
\mathcal L(\theta)
=
\mathbb E_{\mathcal D}
\left[
\frac{1}{N_A}
\sum_{d=1}^{N_A}
\left(
y_t
-
Q_d(s_t,a_{d,t};\theta)
\right)^2
\right].
}
\]
Training uses an experience replay buffer, a periodically or softly updated target
network, and $\epsilon$-greedy exploration.  Numerical learning rate, batch size, replay
capacity, target-update cadence, exploration schedule, total training budget, and seed
set are configuration/training-protocol parameters rather than hard-coded controller
science.

\subsubsection{Stable-Baselines3 implementation boundary}

Stable-Baselines3 2.8.0 provides DQN only for a single \texttt{Discrete} action space and
does not implement the branching Dueling Double DQN algorithm or native DQN support for
the required \texttt{MultiDiscrete}-like action structure~\cite{sb3_280}.  Therefore the
P1 BDQ policy/learning algorithm is implemented explicitly in PyTorch.  Generic
SB3/SB3-style utilities may be reused where they do not alter the ratified BDQ equations
(e.g., replay-buffer patterns, schedules, callbacks, logging, and evaluation utilities).
The implementation must not wrap the full joint action into one enormous flat
\texttt{Discrete} space.

\subsection{Fair-comparison contract}
\label{subsec:p1_fairness}

Within a P1 experiment, RBC, MPC, and RL use:
\begin{itemize}
    \item the same feeder and episode time axis;
    \item the same realized load and available-PV trajectories;
    \item the same selected physical controllable assets;
    \item the same physical device limits;
    \item the same voltage target and reported physical metrics;
    \item the same initial realized device states;
    \item the same control interval;
    \item the same authoritative OpenDSSSim plant.
\end{itemize}
MPC and RL receive the same authorized forecast tensor.  RBC is intentionally defined as a
current-state local heuristic and does not use future forecasts.

The MPC has continuous inverter-$Q$ decisions while BDQ uses the five-level normalized
discretization above.  This is a transparent paradigm-specific action representation, not
a different physical inverter capability.

\subsection{Capability and unsupported-circuit contract}
\label{subsec:p1_capability}

Generalized P1 source must instantiate from circuit/ProblemSpec metadata, not hard-coded
IEEE-123 identifiers.  The required capabilities include:
\begin{itemize}
    \item active bus--phase voltage observations;
    \item selected regulator/tap channels with states, bounds, and command semantics;
    \item selected capacitor channels with states and reactive ratings;
    \item selected PV-inverter channels with active power, apparent-power ratings, and
          reactive command semantics;
    \item load and PV current/exogenous data;
    \item forecast view for MPC/RL;
    \item for MPC, a radial phase-aware non-mutating model service sufficient to construct
          LinDistFlow coefficients/mappings.
\end{itemize}
A controller that cannot instantiate its declared requirements returns an explicit
machine-readable \texttt{UNSUPPORTED} result.  It does not silently degrade to another
physical problem or bypass the plant.

\subsection{Implementation and experiment freeze boundary}
\label{subsec:p1_remaining}

The controller mathematics in this document are ratified.  The following remain
configuration/qualification items before the final PESGM P1 campaign freeze:
\begin{enumerate}
    \item exact IEEE-123 regulator, capacitor, PV device placements, phase mappings, and
          ratings;
    \item the single fixed numerical tuple $(w_V,w_{\mathrm{sw}},w_Q)$ after scale
          verification, with no Pareto optimization;
    \item the RL training-budget/hyperparameter protocol, seeds, and checkpoint-selection
          rule;
    \item measured validation of the LinDistFlow coefficients against the authoritative
          OpenDSSSim plant for the selected experiment feeder.
\end{enumerate}
